% Motor: XeLaTeX, fijado en .latexmkrc y en la receta del editor.
\documentclass[aspectratio=169,10pt,c]{beamer}

% !TEX root = slides.tex
% Preambulo comun: estilo UTP, colores, fuentes, plantillas de beamer y
% comandos de apoyo didactico. Se carga desde slides.tex tras \documentclass.

% La opción "c" de beamer reparte el espacio libre 1:1.5 (arriba:abajo), por lo
% que el contenido queda por encima del centro. Igualamos ambos rellenos.
\makeatletter
\define@key{beamerframe}{c}[true]{%
  \beamer@frametopskip=0pt plus 1fill\relax%
  \beamer@framebottomskip=0pt plus 1fill\relax%
  \beamer@frametopskipautobreak=0pt plus .5\paperheight\relax%
  \beamer@framebottomskipautobreak=0pt plus .5\paperheight\relax%
}
\makeatother

\usepackage{graphicx}
\usepackage{tikz}
\usepackage{xcolor}
\usepackage{iftex}

\ifPDFTeX
  \usepackage[utf8]{inputenc}
  \usepackage[T1]{fontenc}
  \usepackage{lmodern}          % tipografia escalable para el respaldo con pdflatex
  \renewcommand{\familydefault}{\sfdefault}
\else
  \usepackage{fontspec}
  \IfFontExistsTF{Arial}
    {\setsansfont{Arial}[BoldFont={Arial Bold}]}
    {\setsansfont{Helvetica}[BoldFont={Helvetica Bold}]}
\fi

\usepackage[spanish,es-nodecimaldot]{babel}
\usepackage{microtype}
\usepackage{booktabs}
\usepackage{amssymb}
\usepackage{array}
\usetikzlibrary{arrows.meta,positioning,fit,calc,shapes.geometric}

\graphicspath{{assets/}}

% ---------- Identidad visual UTP ----------
\definecolor{utpRojo}{HTML}{C42B35}
\definecolor{utpNegro}{HTML}{111111}
\definecolor{utpGris}{HTML}{5F5F5F}
\definecolor{utpGrisMedio}{HTML}{9A9A9A}
\definecolor{utpGrisClaro}{HTML}{F1F1F1}
\definecolor{utpNaranja}{HTML}{C1670B}

\usetheme{default}
\setbeamercolor{normal text}{fg=utpNegro,bg=white}
\setbeamercolor{structure}{fg=utpRojo}
\setbeamercolor{frametitle}{fg=utpNegro,bg=white}
\setbeamercolor{title}{fg=utpNegro}
\setbeamercolor{subtitle}{fg=utpNegro}
\setbeamercolor{author}{fg=utpNegro}
\setbeamercolor{date}{fg=utpNegro}
\setbeamercolor{block title}{fg=white,bg=utpNegro}
\setbeamercolor{block body}{fg=utpNegro,bg=utpGrisClaro}
\setbeamercolor{alertblock title}{fg=white,bg=utpRojo}
\setbeamercolor{alertblock body}{fg=utpNegro,bg=utpRojo!7}
\setbeamercolor{exampleblock title}{fg=white,bg=utpGris}
\setbeamercolor{exampleblock body}{fg=utpNegro,bg=utpGris!8}

\setbeamerfont{title}{series=\bfseries,size=\fontsize{19.5}{22.5}\selectfont}
\setbeamerfont{subtitle}{size=\fontsize{11.5}{13.5}\selectfont}
\setbeamerfont{date}{series=\bfseries,size=\fontsize{11.5}{13.5}\selectfont}
\setbeamerfont{author}{size=\fontsize{10.5}{12.5}\selectfont}
\setbeamerfont{frametitle}{series=\bfseries,size=\large}
\setbeamerfont{block title}{series=\bfseries,size=\small}
\setbeamerfont{block body}{size=\normalsize}

\setbeamertemplate{navigation symbols}{}
\setbeamertemplate{headline}{}
\setbeamertemplate{blocks}[default]
\setbeamertemplate{itemize item}{\color{utpRojo}\small\raise.15ex\hbox{$\blacktriangleright$}}
\setbeamertemplate{itemize subitem}{\color{utpGris}\scriptsize$\bullet$}
\setbeamertemplate{enumerate item}{\color{utpRojo}\insertenumlabel.}
\setbeamertemplate{background canvas}{\color{white}\rule{\paperwidth}{\paperheight}}

% Título de diapositiva: negrita negra a la izquierda, logo UTP a la derecha
\setbeamertemplate{frametitle}{%
  \begin{tikzpicture}[remember picture,overlay]
    \node[anchor=north east,inner sep=0pt]
      at ([xshift=-7mm,yshift=-5.5mm]current page.north east)
      {\includegraphics[width=17mm]{utp-logo.png}};
  \end{tikzpicture}%
  \vspace{1.2mm}%
  \begin{minipage}[t]{.80\textwidth}\raggedright\insertframetitle\end{minipage}\par
  \vspace{1.1mm}{\color{utpRojo}\rule{\textwidth}{1.1pt}}
}

% Solo numeración discreta; la plantilla UTP no lleva pie de página
\setbeamertemplate{footline}{%
  \hfill{\color{utpGrisMedio}\scriptsize\insertframenumber}\hspace{6mm}\vspace{2.5mm}%
}

% Portada UTP
\setbeamertemplate{title page}{%
  \begin{tikzpicture}[remember picture,overlay]
    \node[anchor=center,inner sep=0pt] at (current page.center) {%
      \includegraphics[width=\paperwidth,height=\paperheight]{utp-cover-bg.jpeg}%
    };
    \node[anchor=north,align=center,text=utpNegro,inner sep=0pt,
      text width=145mm,font=\usebeamerfont{title}]
      at ([yshift=-25.2mm]current page.north) {\inserttitle};
    \node[anchor=north,align=center,text=utpNegro,inner sep=0pt,
      text width=142mm,font=\usebeamerfont{subtitle}]
      at ([yshift=-38.2mm]current page.north) {\insertsubtitle};
    \node[anchor=north,align=center,text=utpNegro,inner sep=0pt,
      font=\usebeamerfont{date}]
      at ([yshift=-50.4mm]current page.north) {\insertdate};
    \node[anchor=north,align=center,text=utpNegro,inner sep=0pt,
      font=\usebeamerfont{author}]
      at ([yshift=-58.4mm]current page.north) {\insertauthor};
  \end{tikzpicture}%
}

% ---------- Comandos de apoyo didáctico ----------
\newcommand{\concepto}[2]{%
  \begin{block}{#1}\large #2\end{block}}
\newcommand{\pregunta}[1]{%
  \begin{alertblock}{Decide y explica}#1\end{alertblock}}
\newcommand{\analogia}[1]{%
  \begin{exampleblock}{En palabras simples}\small #1\end{exampleblock}}
\newcommand{\enlarealidad}[1]{%
  \begin{block}{Cómo se ve en la realidad}\small #1\end{block}}
\newcommand{\ojo}[1]{%
  \begin{alertblock}{Ojo con esto}\small #1\end{alertblock}}
\newcommand{\apoyo}{{\normalfont\footnotesize\color{utpGris}~· material de apoyo}}
\newcommand{\marca}[1]{\tikz[baseline=(m.base)]{%
  \node[rounded corners=1.5pt,fill=utpGrisClaro,draw=utpGrisMedio,inner xsep=3pt,inner ysep=1.5pt,
        font=\footnotesize] (m) {#1};}}

\tikzset{
  caja/.style={draw=utpNegro, very thick, rounded corners=1.5pt, fill=white,
    align=center, minimum height=0.72cm, inner sep=5pt, font=\small},
  cajaazul/.style={caja, draw=utpGris, fill=white},
  cajaverde/.style={caja, draw=utpNegro, fill=utpGrisClaro},
  cajanaranja/.style={caja, draw=utpRojo!75, fill=utpRojo!7},
  flecha/.style={-{Stealth[length=2.3mm]}, very thick, draw=utpNegro},
  flechavuelta/.style={{Stealth[length=2.3mm]}-{Stealth[length=2.3mm]}, very thick, draw=utpNegro}
}

% ---------- Notación pata de gallo (crow's foot) ----------
% La rotación indica hacia dónde está la entidad: 0 derecha, 180 izquierda,
% 90 arriba, -90 abajo. La pata se abre hacia la entidad del lado "muchos".
\tikzset{
  relacion/.style={thick, draw=utpNegro},
  pics/uno/.style={code={%
    \draw[thick,utpNegro,line cap=round] (0,-1.5mm) -- (0,1.5mm);}},
  pics/muchos/.style={code={%
    \draw[thick,utpNegro,line cap=round,line join=round]
      (2.6mm,1.5mm) -- (0,0) -- (2.6mm,-1.5mm);
    \draw[thick,utpNegro,line cap=round] (0,0) -- (2.6mm,0);}},
}
\newcommand{\simbuno}{\tikz[baseline=-0.6ex]{\draw[relacion](0,0)--(7mm,0);\pic at (4.4mm,0){uno};}}
\newcommand{\simbmuchos}{\tikz[baseline=-0.6ex]{\draw[relacion](0,0)--(7mm,0);\pic at (4.4mm,0){muchos};}}


\title{Primeros pasos en SQL}
\subtitle{Del modelo E-R a SQL:\\ nombres de tablas, columnas y tipos de datos}
\date{Semana 06 - Sesión 11}
\author{Rel Guzman}

% Espacio en blanco para completar
\newcommand{\bl}[1][1.9cm]{\underline{\hspace{#1}}}
% Nombre de tabla, siempre en naranja
\newcommand{\tbl}[1]{\textcolor{utpNaranja}{\texttt{#1}}}
% Palabra reservada de SQL dentro de un bloque de código
\newcommand{\kw}[1]{\textcolor{utpRojo}{\textbf{#1}}}
% Parte nueva de un bloque de codigo que crece paso a paso
\newcommand{\nuevo}[1]{{\setlength{\fboxsep}{1.2pt}\colorbox{utpRojo!14}{#1}}}
% Marcas de clave
\newcommand{\pk}{{\scriptsize\color{utpRojo}\textbf{PK}}}
\newcommand{\fk}{{\scriptsize\color{utpNaranja}\textbf{FK}}}
% Insignia de práctica
\newcommand{\ins}[1]{\tikz[baseline=-0.5ex]{\node[rounded corners=2pt,fill=utpRojo,text=white,
  font=\bfseries\scriptsize,inner xsep=5pt,inner ysep=2.5pt] {#1};}}
% Bloque de código SQL sobre fondo gris
\newsavebox{\sqlbuf}
\newenvironment{sqlcode}[1][.86\textwidth]
  {\begin{center}\begin{lrbox}{\sqlbuf}\begin{minipage}{#1}\ttfamily\footnotesize\vspace{6pt}\raggedright}
  {\vspace{6pt}\end{minipage}\end{lrbox}\colorbox{utpGrisClaro}{\usebox{\sqlbuf}}\end{center}}
% Sangría dentro de un bloque de código
\newcommand{\ind}{\hspace*{1.6em}}
% Tabla de tres columnas para alinear nombre / tipo / restricciones
\newenvironment{sqltab}
  {\begin{tabular}{@{}l@{\hspace{1.4em}}l@{\hspace{1.4em}}l@{}}}
  {\end{tabular}}
\newcommand{\sqlline}[1]{\multicolumn{3}{@{}l}{#1}}

\begin{document}

% 1 ---------------------------------------------------------------------
{\setbeamertemplate{footline}{}\maketitle}

% 2 ---------------------------------------------------------------------
\begin{frame}[c]
\begin{block}{Logro de la sesión}
Al finalizar, podrás tomar una entidad del \textbf{dibujo} de una base de datos y escribir una tabla en \textbf{SQL}, con sus columnas y tipos de datos.
\end{block}
\end{frame}

% 3 ---------------------------------------------------------------------
\begin{frame}{Dónde estamos}
\begin{center}
\begin{tikzpicture}
  \node[caja,text width=3.1cm,font=\small,align=center] (a) {\textbf{El dibujo}\\[2pt] \scriptsize el modelo E-R};
  \node[cajanaranja,text width=3.1cm,font=\small,align=center,right=16mm of a] (b) {\textbf{SQL}\\[2pt] \scriptsize lo que se escribe};
  \node[cajaverde,text width=3.1cm,font=\small,align=center,right=16mm of b] (c) {\textbf{La base de datos}\\[2pt] \scriptsize tablas de verdad};
  \draw[flecha] (a)--(b);
  \draw[flecha] (b)--(c);
  \node[font=\scriptsize,below=2mm of a] {sesiones 06 a 09};
  \node[font=\scriptsize,below=2mm of b,text=utpRojo] {\textbf{hoy}};
  \node[font=\scriptsize,below=2mm of c] {a partir de hoy};
\end{tikzpicture}
\end{center}
\vspace{0.55cm}
\begin{block}{La idea de la sesión}
\small El dibujo ya está hecho. Falta \textbf{escribirlo}, porque nadie puede guardar un pedido dentro de un dibujo.
\end{block}
\end{frame}

% 4 ---------------------------------------------------------------------
\begin{frame}[plain]
  \begin{tikzpicture}[remember picture,overlay]
    \node[anchor=north east,inner sep=0pt]
      at ([xshift=-10mm,yshift=-8mm]current page.north east)
      {\includegraphics[width=43mm]{utp-logo.png}};
    \node[anchor=center,inner sep=0pt] at (current page.center)
      {\resizebox{70mm}{!}{\bfseries\color{utpNegro}¿Qué es SQL?}};
    \node[anchor=north,inner sep=0pt] at ([yshift=-11mm]current page.center)
      {\color{utpRojo}\rule{70mm}{1.6pt}};
  \end{tikzpicture}
\end{frame}

% 5 ---------------------------------------------------------------------
\begin{frame}{¿Qué es SQL y para qué sirve?}
\concepto{SQL}{\small El idioma con el que se le \textbf{pide} algo a una base de datos.}
\vspace{0.4cm}
\begin{center}
\resizebox{\textwidth}{!}{%
\begin{tikzpicture}
  \node[caja,text width=2.2cm,font=\scriptsize,align=center] (yo) {\textbf{tú}};
  \node[cajanaranja,text width=2.8cm,font=\scriptsize,align=center,right=16mm of yo] (sql) {\textbf{una frase\\ en SQL}};
  \node[caja,text width=2.6cm,font=\scriptsize,align=center,right=16mm of sql] (ges) {\textbf{el gestor}\\ \scriptsize MySQL};
  \node[cajaverde,text width=2.4cm,font=\scriptsize,align=center,right=16mm of ges] (bd) {\textbf{las tablas}};
  \draw[flecha] (yo)--(sql); \draw[flecha] (sql)--(ges); \draw[flecha] (ges)--(bd);
  \node[font=\scriptsize,below=2mm of sql,text width=3.2cm,align=center] {escribes \textbf{qué} quieres};
  \node[font=\scriptsize,below=2mm of ges,text width=3.2cm,align=center] {él decide \textbf{cómo} hacerlo};
\end{tikzpicture}}
\end{center}
\vspace{0.45cm}
\small Sirve para \textbf{crear} las tablas (hoy), \textbf{guardar} filas y \textbf{consultarlas} (las próximas sesiones). Es un \textbf{estándar}: lo mismo sirve en MySQL, PostgreSQL, SQL Server u Oracle.
\end{frame}

% 6 ---------------------------------------------------------------------
\begin{frame}{¿Y dónde se escribe SQL?}
\small En un programa que se conecta al gestor. Todos hacen lo mismo:
\begin{itemize}\setlength\itemsep{3pt}\small
  \item MySQL Workbench
  \item DBeaver
  \item phpMyAdmin \textnormal{(en el navegador)}
  \item la consola \texttt{mysql} \textnormal{(solo texto)}
\end{itemize}
\vspace{6pt}
\small Escribes la instrucción, la ejecutas, y el gestor \textbf{responde}.
\end{frame}

% 7 ---------------------------------------------------------------------
\begin{frame}{Una instrucción se lee casi como una frase}
\begin{sqlcode}[.64\textwidth]
\kw{CREATE TABLE} cliente ( \ldots );
\end{sqlcode}
\vspace{0.2cm}
\begin{center}
\small
\begin{tabular}{>{\ttfamily}l>{\raggedright\arraybackslash}p{6.4cm}}
\toprule
\normalfont\bfseries En SQL & \bfseries En español \\
\midrule
CREATE & crea \\
TABLE & una tabla \\
cliente & que se llame cliente \\
( \ldots ) & con estas columnas \\
; & y aquí termino de pedirte \\
\bottomrule
\end{tabular}
\end{center}
\vspace{0.4cm}
\begin{block}{Tres costumbres, no reglas}
\small SQL \textbf{no} distingue mayúsculas de minúsculas, pero: las palabras del idioma van en \textbf{MAYÚSCULA}, los nombres que inventamos en minúscula, y cada instrucción termina en \textbf{punto y coma}.
\end{block}
\end{frame}

% 8 ---------------------------------------------------------------------
\begin{frame}{\ins{Práctica 1} Lee y escribe qué hace}
\begin{center}
\small
\begin{tabular}{>{\ttfamily}p{5.4cm}>{\raggedright\arraybackslash}p{6.2cm}}
\toprule
\normalfont\bfseries La instrucción & \bfseries Qué le estoy pidiendo al gestor \\
\midrule
CREATE DATABASE colegio; & \bl[5.4cm] \\[9pt]
CREATE TABLE alumno ( \ldots ); & \bl[5.4cm] \\[9pt]
DROP TABLE alumno; & \bl[5.4cm] \\[9pt]
USE colegio; & \bl[5.4cm] \\
\bottomrule
\end{tabular}
\end{center}
\vspace{0.4cm}
\small \textbf{Pista:} traduce palabra por palabra. \texttt{CREATE} es «crea», \texttt{DROP} es «elimina», \texttt{USE} es «usa». Lo que sigue dice \textbf{qué cosa}.
\end{frame}

% 9 ---------------------------------------------------------------------
\begin{frame}{Paso 1: el nombre de la tabla}
\begin{columns}[c,onlytextwidth]
\column{.56\textwidth}
\begin{sqlcode}[\textwidth]
\nuevo{\kw{CREATE TABLE} cliente}\\
(\\
);
\end{sqlcode}
\column{.40\textwidth}
\begin{center}
\begin{tikzpicture}
  \node[cajaverde,text width=2.7cm,font=\scriptsize,align=left]
    {\textbf{cliente}\\[3pt] \texttt{id} \ \pk\\ \texttt{nombre}\\ \texttt{correo}};
\end{tikzpicture}
\end{center}
\vspace{4pt}
\centering\scriptsize vamos a escribir\\ este dibujo
\end{columns}
\vspace{0.5cm}
\small \texttt{CREATE TABLE} es «crea una tabla nueva», y \texttt{cliente} es cómo se va a llamar. Entre los \textbf{paréntesis} van a ir las columnas: por ahora está vacío.
\vspace{0.2cm}

\small El nombre no se inventa: sale del dibujo, donde la entidad ya se llamaba \tbl{cliente}.
\end{frame}

% 10 ---------------------------------------------------------------------
\begin{frame}{Paso 2: la primera columna}
\begin{sqlcode}[.56\textwidth]
\kw{CREATE TABLE} cliente\\
(\\
\ind \nuevo{nombre \ VARCHAR(80)}\\
);
\end{sqlcode}
\vspace{0.35cm}
\begin{center}
\begin{tikzpicture}
  \node[cajaverde,text width=2.4cm,font=\scriptsize,align=center] (n) {\texttt{nombre}};
  \node[font=\scriptsize,below=2mm of n] {cómo se llama};
  \node[cajanaranja,text width=2.9cm,font=\scriptsize,align=center,right=18mm of n] (t) {\texttt{VARCHAR(80)}};
  \node[font=\scriptsize,below=2mm of t] {qué se puede guardar};
\end{tikzpicture}
\end{center}
\vspace{0.55cm}
\small \textbf{Una columna son siempre dos cosas: un nombre y un tipo.} Aquí estamos diciendo «guarda texto, de hasta 80 caracteres».
\end{frame}

% 11 ---------------------------------------------------------------------
\begin{frame}{Paso 3: las demás columnas}
\begin{sqlcode}[.56\textwidth]
\kw{CREATE TABLE} cliente\\
(\\
\ind nombre \ VARCHAR(80)\nuevo{,}\\
\ind \nuevo{correo \ VARCHAR(120)}\\
);
\end{sqlcode}
\vspace{0.35cm}
\small Se agregan una debajo de otra, y \textbf{la coma las separa}.
\vspace{0.4cm}
\begin{alertblock}{El error número uno del primer día}
\small La coma \textbf{separa} columnas, no las termina: \textbf{la última no lleva coma}. Si sobra una, el gestor responde \texttt{You have an error in your SQL syntax\ldots near ')'} y señala la línea \textbf{siguiente} a la del problema.
\end{alertblock}
\end{frame}

% 12 ---------------------------------------------------------------------
\begin{frame}[plain,c]
\centering
{\fontsize{30}{36}\selectfont\bfseries
Resolvemos un poco\\[6pt]
de la práctica\par}
\end{frame}

\end{document}
